\section{Distributed maintenance as a network of maintenance relations}

The preceding sections suggest a representation richer than ``controller acts on state.'' \Cref{fig:maintenance} summarizes the proposed decomposition.

\begin{figure}[htbp]
\centering
\begin{tikzpicture}[
  font=\small,
  >=Latex,
  box/.style={draw, rounded corners, align=center, minimum height=10mm, text width=28mm, inner sep=3pt},
  wide/.style={draw, rounded corners, align=center, minimum height=10mm, text width=34mm, inner sep=3pt},
  note/.style={draw, dashed, rounded corners, align=center, text width=34mm, inner sep=3pt},
  arr/.style={->, thick}
]
\node[wide] (V) at (0,3.4) {Declared viability condition \(\V\)\\and disturbance class \(\D\)};
\node[box] (S) at (4,3.4) {Shared enabling\\state};
\node[box] (F) at (8,3.4) {Viability-relevant\\local function \(F_i\)};
\node[box] (Y) at (8,1.15) {Locally available\\signal / proxy \(Y_i\)};
\node[box] (R) at (4,1.15) {Return / renewal\\\(R_i=G(Y_i)\)};
\node[box] (C) at (0,1.15) {Future corrective\\capacity \(C_i\)};

\draw[arr] (V) -- node[above]{defines} (S);
\draw[arr] (S) -- (F);
\draw[arr] (F) -- (Y);
\draw[arr] (Y) -- (R);
\draw[arr] (R) -- (C);
\draw[arr] (C.north) |- (S.west);

\node[note] (D) at (0,-1.65) {\textbf{Failure domains}\\Nominal multiplicity need not supply independent correction};
\node[note] (M) at (4,-1.65) {\textbf{Multiple margins}\\\(\A=\bigcap_j\A_j\) may be broad, narrow, or empty};
\node[note] (P) at (8,-1.65) {\textbf{Proxy gap}\\\(Y_i\) may remain high while viability contribution \(F_i\) falls};

\draw[arr,dashed] (D.north) -- (C.south);
\draw[arr,dashed] (P.north) -- (Y.south);

\node[fit=(S)(F)(Y)(R)(C), draw, densely dotted, rounded corners, inner sep=6pt] {};
\end{tikzpicture}

\caption{\textbf{Distributed maintenance as a recursive viability architecture.}
A declared viability condition specifies what continuation means under a disturbance class. Local processes receive function-relevant information, act on shared state, and depend on return or renewal for future corrective capacity. Failure-domain structure determines whether nominal multiplicity supplies independent correction. The function--renewal proxy gap appears when the signal that attracts return diverges from the function relevant to viability. Multiple viability conditions generate intersecting admissible regions rather than a single scalar objective.}
\label{fig:maintenance}
\end{figure}

For each local process, distinguish at least the following relations.

\subsection{Viability specification}

The continuation property and disturbance class must be stated. Without this step, ``maintenance'' has no definite target.

\subsection{Informative coupling}

Some local variable must carry information relevant to the condition being maintained. This need not mean representation in a cognitive sense.

\subsection{Corrective or adaptive effect}

Changes in the local process must alter future viability. Otherwise the signal is epiphenomenal with respect to the declared condition.

\subsection{Independent corrective capacity}

Where common-mode failure is possible, nominal multiplicity is not equivalent to independent correction. Reliability depends on failure topology.

\subsection{Return-loop closure}

Processes that consume resources or decay must receive sufficient renewal to continue supplying their function. The return may be metabolic, energetic, structural, reproductive, economic or informationally mediated.

\subsection{Renewal of maintenance capacity}

The ability to correct future disturbance must itself survive quiet periods, use, degradation and replacement.

\subsection{Architectural corrigibility}

The maintenance architecture can itself become maladaptive. A system able only to correct its object-level state, but unable to revise a failing maintenance mechanism, remains vulnerable to structural change outside the mechanism's operating assumptions.

These seven items are proposed as a diagnostic decomposition. They are not claimed to be universally necessary and sufficient conditions for every persistent network.

\section{Worked comparison: engineered auditing and fungal transport}

The purpose of comparing radically different substrates is not to assert mechanism identity. It is to ask what remains meaningful when obvious substrate-specific concepts are removed.

\subsection{Engineered distributed verification}

In a distributed verification protocol, one can identify a validity condition, randomly selected verification processes, explicit judgments, escalation after detectable failure, implementation diversity, fault records, and externally supplied economic return.

The architecture is highly explicit. Signals, actions and participants have designed semantics. Yet the analysis reveals that nominally honest participation is not automatically equivalent to independent corrective capacity, and that an observable output such as a judgment need not by itself establish the costly process presumed to generate it.

The companion JAM/ELVES manuscript develops those claims at protocol level. They are not repeated here as evidence of a generic theorem.

\subsection{Fungal transport}

A fungal network contains no validator, voting rule or staking contract. Its transport architecture nevertheless involves flows through network paths, local structural growth and regression, redistribution under changing resource conditions, topology-dependent transport, and trade-offs among connectivity, transport efficiency, robustness and construction cost.

Empirical work shows that fungal network traits occupy such trade-offs rather than maximizing a single objective \citep{aguilar2022}.

The useful comparison is therefore not
\[
\text{validator}=\text{fungal cord},
\]
nor
\[
\text{reward}=\text{nutrient}.
\]

The defensible cross-substrate statement is weaker:

\begin{quote}
In both systems, future allocation to a local process depends on information generated by its interaction with a wider network, and that allocation alters future network capacity.
\end{quote}

The implementation differs radically. The maintenance relation survives.

\subsection{What the comparison removes}

The comparison shows why several terms should not be primitives of the general framework: \emph{agent} is too restrictive; \emph{monitor} is too architectural; \emph{reward} is too economic; \emph{commons} is too domain-specific; and \emph{individual} may be ambiguous across scale.

More neutral terms are
\[
\boxed{
\text{local process}
\rightarrow
\text{shared enabling condition}
\rightarrow
\text{return to future process capacity}.
}
\]

This abstraction does not imply that biological and engineered systems belong to one causal category. It is a common bookkeeping language for asking comparable viability questions.

\section{Relationship to organizational autonomy}

The resulting description has obvious proximity to organizational theories of autonomy.

Mossio and Moreno characterize organisms by organizational closure: constraints depend on other constraints for their own production while contributing to processes that maintain other constraints \citep{mossio2010}. Montévil and Mossio formalize biological organization as closure of constraints \citep{montevil2015}. Bich and colleagues argue that regulation is not exhausted by constitutive closure: regulatory mechanisms sense relevant conditions and modulate other constraints, amounting to second-order control \citep{bich2016}. More recent work emphasizes heterarchical coordination of such mechanisms rather than requiring a central regulator \citep{bich2022}.

The present framework shares three commitments with this tradition:
\begin{enumerate}[label=(\arabic*)]
\item functions are relational rather than intrinsic;
\item maintenance processes are themselves dependent on other processes;
\item regulation can operate through distributed control of control.
\end{enumerate}

It differs mainly in scope and emphasis.

First, biological autonomy is not assumed. Engineered networks are included even when their energy, replacement, governance and reward loops lie partly outside the focal system.

Second, the framework begins with a \emph{declared viability criterion} rather than claiming that organizational closure alone determines the relevant function.

Third, contribution remains explicitly vector-valued and architecture-dependent.

Fourth, the framework puts special emphasis on observable maintenance diagnostics: current state, corrective capacity, independent failure domains, renewal pathways and function--proxy mismatch.

Distributed maintenance is therefore best understood as a complementary analytical layer, not a competitor to organizational autonomy.
